% ClassicThesis 4.1, pinned from the repository requested by the author.
\RequirePackage{fix-cm}
\documentclass[oneside,openany,titlepage,numbers=noenddot,headinclude,
  footinclude=true,cleardoublepage=empty,BCOR=0mm,paper=a4,fontsize=11pt]{scrbook}
% Project configuration, based on the upstream classicthesis-config.tex.
% Typography and headings belong to ClassicThesis; the digital page geometry,
% running heads and code listings are configured here, after upstream setup.
% The vendored upstream package is deliberately not modified.
\PassOptionsToPackage{dvipsnames}{xcolor}
\usepackage{fontspec}
\usepackage{polyglossia}
\setdefaultlanguage{slovak}
\usepackage{amsmath,amssymb}
\usepackage{booktabs,longtable,array,tabularx}
\usepackage{graphicx}
\usepackage{caption}
\captionsetup{format=hang,font=small,labelfont=normalfont}
\usepackage{listings}
\usepackage{enumitem}
\usepackage{needspace}
\usepackage{tikz}
\usetikzlibrary{arrows.meta,positioning,calc}
\usepackage[backend=biber,style=numeric,sorting=none,giveninits=true,maxbibnames=6]{biblatex}
\addbibresource{references.bib}
\usepackage{hyperref}
\usepackage{xurl}
\usepackage[eulerchapternumbers,floatperchapter,manychapters]{classicthesis}
% ClassicThesis first establishes its typearea. Geometry is the final authority:
% do not call \areaset, \recalctypearea or \newgeometry later in the document.
\usepackage[a4paper,left=30mm,right=30mm,top=23mm,bottom=25mm,
  includeheadfoot,headheight=15pt,headsep=14pt,footskip=26pt,
  bindingoffset=0pt]{geometry}
% A stable chapter title is more useful here than the repeated section name
% "Prednáška". All pages use the same right-aligned head and a restrained rule.
\clearpairofpagestyles
\KOMAoptions{headwidth=text,footwidth=text,headsepline=0.35pt}
\automark{chapter}
\renewcommand{\chaptermark}[1]{\markboth{#1}{#1}}
\renewcommand{\sectionmark}[1]{}
\setkomafont{pagehead}{\normalfont\small}
\ohead{\normalfont\small\leftmark}
\cfoot[\pagemark]{\pagemark}
% Unicode equivalents of ClassicThesis' original Palatino family.
\setmainfont{TeX Gyre Pagella}[Numbers=OldStyle]
\setsansfont{TeX Gyre Heros}
\setmonofont{DejaVu Sans Mono}[Scale=MatchLowercase]
\input{unicode-spacing}
% Retain the upstream chapter font, title, rule and spacing. Only place the
% Euler numeral inside the text block: "10" would overrun a 30mm side margin.
\titleformat{\chapter}[display]{\relax}
  {\makebox[\linewidth][r]{\smash{\color{halfgray}\chapterNumber\thechapter}}}{0pt}
  {\raggedright\spacedallcaps}[\normalsize\vspace*{.8\baselineskip}\titlerule]
% Distinguish topic headings from both prose and italic technical vocabulary.
% Keep the upstream chapter treatment; give smaller headings explicit weight.
\titleformat{\section}{\normalfont\Large\bfseries}{\thesection}{0.7em}{}
\titleformat{\subsection}{\normalfont\large\bfseries}{\thesubsection}{0.7em}{}
\titleformat{\subsubsection}{\normalfont\normalsize\bfseries}{\thesubsubsection}{0.7em}{}
\titlespacing*{\section}{0pt}{1.6\baselineskip}{0.65\baselineskip}
\titlespacing*{\subsection}{0pt}{1.3\baselineskip}{0.5\baselineskip}
\titlespacing*{\subsubsection}{0pt}{1\baselineskip}{0.4\baselineskip}
% Paragraphs use one consistent convention, including after headings/lists.
\setlength{\parindent}{0pt}
\setlength{\parskip}{0.45\baselineskip plus 1pt minus 0.5pt}
% Color is confined to the cover and restrained syntax highlighting.
\colorlet{Maroon}{black}
\definecolor{muted}{gray}{0.4}
\hypersetup{unicode=true,hidelinks,bookmarksnumbered=true,
  pdftitle={@@TITLE@@},pdfauthor={@@AUTHOR@@},pdfsubject={@@SUBTITLE@@}}
\setcounter{tocdepth}{1}
\setcounter{secnumdepth}{1}
\setlist{topsep=0.4\baselineskip,itemsep=0.15\baselineskip,parsep=0pt,leftmargin=*}
\setlength{\emergencystretch}{3em}
\clubpenalty=10000
\widowpenalty=10000
% One style for presentation fragments, worksheets, Markdown and external files.
% Native listings remain page-breakable: no float, minipage or enclosing box.
\definecolor{codebackground}{gray}{0.965}
\definecolor{codekeyword}{RGB}{35,61,83}
\definecolor{codestring}{RGB}{103,75,55}
\definecolor{codecomment}{gray}{0.40}
\definecolor{codenumber}{gray}{0.48}
\lstdefinelanguage{bookjson}{
  morestring=[b]",
  morekeywords={true,false,null},
  sensitive=true
}
% listings' default character table ends at U+00FF. Under XeTeX, characters
% above that range otherwise bypass its token buffer and can move within words
% (e.g. "Poslať" becomes "ťPosla"). Extend the table using its own processors,
% locally for each listing; retain the literal glyphs, including in strings.
\makeatletter
\begingroup
\catcode0=12\relax
\g@addto@macro\lst@DefEC{%
  \lst@CCECUse\lst@ProcessLetter čďěĺľňŕřšťůžČĎĚĹĽŇŔŘŠŤŮŽ^^00%
  \lst@CCECUse\lst@ProcessOther —€→←─│┌┐└┘├┤┬┴┼^^00%
}
\endgroup
\makeatother
\lstdefinestyle{bookcode}{
  language={},
  basicstyle=\small\ttfamily\color{black},
  keywordstyle=\bfseries\color{codekeyword},
  commentstyle=\itshape\color{codecomment},
  stringstyle=\color{codestring},
  columns=fullflexible,
  keepspaces=true,
  tabsize=4,
  showspaces=false,
  showtabs=false,
  showstringspaces=false,
  upquote=true,
  extendedchars=true,
  numbers=left,
  numberstyle=\scriptsize\ttfamily\color{codenumber},
  numbersep=9pt,
  firstnumber=1,
  stepnumber=1,
  numberblanklines=true,
  breaklines=true,
  breakatwhitespace=false,
  breakautoindent=true,
  breakindent=1em,
  backgroundcolor=\color{codebackground},
  frame=single,
  framerule=0pt,
  frameround=ffff,
  framesep=6pt,
  % The frame starts at the surrounding left margin. Its background includes
  % the number gutter, with 6pt padding outside it and beside the last column.
  xleftmargin=30pt,
  xrightmargin=6pt,
  framexleftmargin=24pt,
  framexrightmargin=0pt,
  aboveskip=0.8\baselineskip,
  belowskip=0.8\baselineskip,
  captionpos=t
}
\lstset{style=bookcode}
\lstnewenvironment{Code}[1][]{\lstset{style=bookcode,#1}}{}

\DeclareCaptionType{listing}[Ukážka][Zoznam ukážok]
\captionsetup[listing]{hypcap=false}
\counterwithin{listing}{chapter}
\newcommand{\sourcefile}[1]{\par{\footnotesize\color{muted}Zdroj: \path{#1}}\par}
\newcommand{\editorial}[1]{\par\smallskip{\small\itshape Redakčná poznámka. #1}\par\smallskip}
\newcounter{exercise}[chapter]
\renewcommand{\theexercise}{\thechapter.\arabic{exercise}}
\newcommand{\exerciseheading}[2]{\refstepcounter{exercise}\label{#1}\subsection*{Úloha \theexercise\enspace #2}}
\tikzset{book node/.style={draw=black!65,line width=0.35pt,rectangle,align=center,
  font=\footnotesize,inner xsep=3pt,inner ysep=4pt},
  book edge/.style={-{Stealth[length=3pt,width=3pt]},line width=0.4pt}}

\begin{document}
\frenchspacing
\raggedbottom
\frontmatter
\pagestyle{plain}
% Cover-only composition. No page geometry or body heading changes here.
\begin{titlepage}
  \centering
  \definecolor{coverink}{RGB}{35,61,83}
  \newcommand{\coverseparator}{%
    \begin{tikzpicture}[baseline=0pt,color=coverink,line width=0.4pt]
      \draw (-52mm,0) -- (-3mm,0);
      \draw (3mm,0) -- (52mm,0);
      \draw (0,1.5mm) -- (1.5mm,0) -- (0,-1.5mm) -- (-1.5mm,0) -- cycle;
    \end{tikzpicture}%
  }
  \vspace*{8mm}
  {\large\spacedlowsmallcaps{@@UNIVERSITY@@}\par}
  \vspace{3mm}
  {\normalsize @@FACULTY@@\par}
  \vspace{\stretch{2}}
  \coverseparator\par
  \vspace{12mm}
  {\color{coverink}\fontsize{32}{40}\selectfont
    \parbox{0.94\textwidth}{\centering @@TITLE@@}\par}
  \vspace{8mm}
  {\large\parbox{0.85\textwidth}{\centering @@SUBTITLE@@}\par}
  \vspace{12mm}
  \coverseparator\par
  \vspace{16mm}
  {\large @@AUTHOR@@\par}
  \vspace{\stretch{3}}
  {\normalsize @@PLACE@@, @@YEAR@@\par}
  \vspace{3mm}
  {\small\spacedlowsmallcaps{@@EDITION@@}\par}
  \vspace*{8mm}
\end{titlepage}

\thispagestyle{empty}
\noindent\spacedlowsmallcaps{@@TITLE@@}\par
@@AUTHOR@@\par
@@UNIVERSITY@@\\@@FACULTY@@\par
@@PLACE@@, @@YEAR@@\par
\vfill
Tento súbor je automaticky zostaveným pracovným rukopisom zo zdrojov predmetu.
Nejde o recenzovanú publikáciu. Bibliografické údaje vydavateľa, recenzentov a ISBN
sa doplnia až pri príprave konkrétneho vydania.
\par\smallskip
Zdrojová revízia: \texttt{@@REVISION@@}. Variant: @@VARIANT@@.
\cleardoublepage
\pagestyle{scrheadings}
\addchap{Predhovor}
Tieto skriptá sú určené študentom informatiky, ktorí už poznajú základy programovania.
Python predstavujú ako nástroj na automatizáciu, tvorbu aplikácií, prácu s dátami
a webovými službami. Každá kapitola spája výklad s praktickým cvičením.
\par
Pri štúdiu najprv predpovedajte správanie ukážky, potom ju spustite a zmeňte vstupy.
Úlohy riešte pred čítaním referenčných implementácií. Pri použití AI nástrojov
musíte vedieť výsledný program vysvetliť a overiť jeho správanie.
\par
Kód z prednášok obsahuje aj zámerne chybné ukážky používané pri rozbore problému.
Ich význam určujú nadpisy a sprievodný výklad. Jednotlivé fragmenty nemusia byť
samostatne spustiteľnými programami. Vstupné dáta a štartovacie súbory zostávajú
v adresároch cvičení a odkazuje na ne sprievodný zoznam zdrojov.
\par
@@PREFACE@@
\cleardoublepage
@@AUTHOR_BIO@@
\pdfbookmark[0]{\contentsname}{contents}
\tableofcontents
\mainmatter
@@CHAPTERS@@
\appendix
@@APPENDICES@@
\backmatter
\nocite{*}
\printbibliography[heading=bibintoc,title={Literatúra a dokumentácia}]
\addchap{Poznámka k zostaveniu}
Obsah vzniká z verejných slovenských prezentácií, pracovných listov a lokálnych
autorských podkladov. Generátor nemení vstupné súbory. Každý zahrnutý zdroj má
kontrolný súčet v sprievodnom súbore \texttt{sources.json}. Zmenou zdroja sa pri
ďalšej kompilácii aktualizuje príslušná kapitola.
\par
Text prednášok je prevzatý z výučbových materiálov; hovorené poznámky ešte vyžadujú
redakčnú úpravu do knižného jazyka. Aktuálny stav úplnosti a revízií uvádza
sprievodný súbor \texttt{report.md}.
\end{document}
